\begin{textsample}{POS Dim 1 – ce – Score 51.00 – ce/ce\_em\_48.txt}  \label{ex:f1_pos_046}
8.4.6 A diversidade e a multiplicidade no ensino de Filosofia : narrativas, usos e formas de construção do conhecimento A palavra [ filosofia ] faz referência a \textbf{uma} relação com o saber e, em particular, a \textbf{um} vínculo de \textbf{amor} como aspiração ou desejo de saber, mais \textbf{que} ao domínio de \textbf{um} saber determinado.
Ao assumir essa caracterização [... ], se se trata de curso \textbf{que} se situa na filosofia — isto é, aquele \textbf{que} poderíamos chamar cabalmente “ filosófico ” —, o \textbf{que} aparece como fundante não é tanto o recorte ocasional de \textbf{um} conhecimento a ser transmitido, mas a atividade de aspirar a “ alcançar o saber ” ( CERLETTI, 2009, p. 18-19 ).
Em sentido estrito, o perguntar filosófico não se detém nunca, porque, para \textbf{um} filósofo, o \textbf{amor} ou desejo de saber ( a filo – sofia ) nunca é preenchida.
Nessa inquietude de saber, a pergunta filosófica se dirige, com perseverança, ao coração do conceito ( CERLETTI, 2009, p. 25 ).
Os trechos escritos por Cerletti abrem trilhas para pensarmos o ensino de Filosofia na última etapa da educação básica, na forma de \textbf{disciplina}, práticas ou estudos.
Partindo da raiz semântica do termo filosofia, filo ( \textbf{amor} ou amizade ) e sofia ( saber ), o filósofo é aquele \textbf{que}, no seu anseio por saber, por amá -lo, deixa clara a importância \textbf{que} a sofia tem na sua vida.
Essa busca, presente na relação do \textbf{homem} com o saber, nunca terá \textbf{um} fim.
Essa perceptiva, com raízes no pensamento de Sócrates, é a mesma \textbf{que} os professores de Filosofia do estado do Ceará podem cultivar : por amar o saber e ter ciência do seu valor e importância, precisa -se desenvolver nas alunas/os a capacidade de amar o saber ou o exercitar aquilo \textbf{que} o leva ao “ coração do conceito ”, dando -os condições de vivenciar a cidadania em sua \textbf{integralidade}, pois esta só se realiza na plenitude quando a participação na vida política é consciente, o \textbf{que} é \textbf{fomentado} pelo trabalho \textbf{conceitual} específico da Filosofia.
Desde a constituição de 1988, a relação da filosofia com a cidadania na educação básica foi expressa, legalmente, com a criação da LDB em 1966.
No artigo Art. 36, §1º, estava expresso \textbf{que} todo \textbf{educando} deveria demonstrar “ domínio dos conhecimentos de Filosofia e de Sociologia necessários ao exercício da cidadania ”. \textbf{Contudo}, a ausência de instrumentos para pensar o modo, o tempo e os conteúdos trabalhados nas aulas dessas duas \textbf{disciplinas} pode ter \textbf{fomentado} a desatenção dos estados da federação ao ensino de filosofia, observado na baixa oferta de aulas, pois a maioria ofertava somente \textbf{uma} aula semanal em \textbf{um} dos anos do ensino médio, realidade modificada com a lei 10.648 de 2008, \textbf{que} tornou o ensino de Filosofia obrigatório em toda a etapa final da Educação Básica.
Em 2017, com a aprovação da lei do Novo Ensino Médio ( 13.415/2017 ), ocorrem mudanças significativas \textbf{que} atingem, direta e indiretamente, o ensino de Filosofia.
O termo “ \textbf{disciplina} obrigatória ”, por exemplo, é substituído pela frase : “ obrigatoriamente estudos e práticas de [... ] filosofia ”.
É implementada a ideia de flexibilização curricular na qual o limite de horas aulas dedicadas à BNCC passa a ser de 1600 \textbf{que}, se dividida hipoteticamente pelos três anos, cada \textbf{um} com 200 dias letivos, as escolas ofertaram 15 horas semanais de base⁴.
No contexto da BNCC e do Novo Ensino Médio há \textbf{um} risco latente dos \textbf{discentes} não terem acesso, em sua \textbf{integralidade}, como obrigava a lei 10.648/2008, dos objetos de aprendizagem presentes no estudo da Filosofia, pois \textbf{uma} escola pode determinar \textbf{que} ela será ministrada, por exemplo, em somente \textbf{um} ano e com \textbf{uma} carga horária reduzida.
Na lógica \textbf{que} definiu os direitos de aprendizagem como competências e habilidades e flexibilizou a carga horária, os conteúdos se tornam secundários, haja vista serem meios, não o fim.
Por outro \textbf{lado}, a BNCC, em todas as suas três versões, levanta \textbf{um} debate novo em torno dos objetivos da educação, \textbf{que} desde a constituição de 1988 se apresenta como “ pleno desenvolvimento do \textbf{educando}, seu preparo para o exercício da cidadania e sua qualificação para o trabalho ”.
O ensino de Filosofia deixa de assumir, unicamente, a função de proporcionar o “ domínio dos conhecimentos [... ] necessários ao exercício da cidadania ” e passa a contribuir na formação integral dos \textbf{aprendizes}.
Desse modo, o papel da Filosofia se torna maior, pois \textbf{uma} formação integral \textbf{exige} \textbf{que} sua especificidade, o trabalho com conceitos filosóficos, seja aprendida, o \textbf{que} se dará a partir do trabalho pedagógico com os problemas e os conteúdos de Filosofia.
Segundo a BNCC, deve -se entender por educação integral : > [... ] formação e [... ] desenvolvimento humano global, o \textbf{que} \textbf{implica} compreender a complexidade e a não linearidade desse desenvolvimento, rompendo com visões reducionistas \textbf{que} privilegiam ou a dimensão intelectual ( cognitiva ) ou a dimensão afetiva.
Significa, ainda, assumir \textbf{uma} visão plural, singular e integral da criança, do adolescente, do jovem e do adulto – considerando -os como \textbf{sujeitos} de aprendizagem – e promover \textbf{uma} educação voltada ao seu acolhimento, reconhecimento e desenvolvimento pleno, nas suas \textbf{singularidades} e diversidades.
Além disso, a escola, como espaço de aprendizagem e de democracia inclusiva, deve se fortalecer na prática coercitiva de não discriminação, não preconceito e respeito às diferenças e diversidades. ( BNCC, 2017, p. 14 ).
Essa ideia de educação integral \textbf{necessita} das diversas narrativas filosóficas, pois proporcionarão aos \textbf{discentes} compreender o significado de juventude, complexidade do mundo, \textbf{integralidade} do ser, sentido da vida, aprender a aprender, cidadania, trabalho, diversidade, diálogo, corpo, sentimentos, linguagem, ciência, arte, \textbf{métodos} e tecnologia, além dos problemas sociais, econômicos, políticos, éticos.
A Filosofia também potencializa o estudo de conceitos mais presentes nas ciências humanas, como abstração, cultura, simbolização, diálogo, \textbf{método} científico ( hipóteses e argumentos ), dúvida, protagonismo, tempo, espaço, temporalidade, espiritualidade, \textbf{territorialidade}, natureza, sociedade, indivíduo, coletividade, dentre outros.
Assim, desde o início da atividade humana, em torno dos seus problemas e questões, a Filosofia é \textbf{um} saber \textbf{basilar} para compreender a multiplicidade da vida, \textbf{visto} no fato de \textbf{que} grande número de \textbf{cientistas}, historiadores, geógrafos, sociólogos, matemáticos, linguistas, lançam dela mão para compreender seus problemas.
Retomando a leitura \textbf{moderna} de Descartes sobre as ciências, a Filosofia seria a mãe das outras ciências.
Algo parecido \textbf{diz} Hegel, quando afirma \textbf{que} a Filosofia \textbf{enquanto} conceito seria, \textbf{enquanto} ciência, a única \textbf{que} pensa as demais ciências, \textbf{enquanto} elas só refletem a si mesmas.
Nesse âmbito se pode \textbf{argumentar} como as diversas narrativas filosóficas, como a Antropologia, a metafísica, a lógica, a ética, a estética, a política, a filosofia das ciências e a filosofia da linguagem, são \textbf{basilares} na formação de qualquer \textbf{sujeito} no presente, principalmente os educandos brasileiros \textbf{que} passam pelo Ensino Médio.
Para \textbf{uma} educação ou \textbf{uma} formação integral \textbf{que} proporcione a igualdade, o fomento à vivência das diversidades e a \textbf{equidade} aos aprendentes, como \textbf{diz} a constituição, a LDB, as Diretrizes Curriculares Nacionais da Educação Básica e a BNCC, o ensino de Filosofia é \textbf{imprescindível}.
Mas há \textbf{um} ensino de Filosofia adequado segundo a DCRD?
Essa resposta, como em todos os componentes curriculares do Ensino Médio, não pode ser unívoca.
Quanto à Filosofia, retomando a citação inicial de Alejandro Cerletti, o docente deve \textbf{ver} sua prática como \textbf{uma} relação entre o \textbf{sujeito} e o saber \textbf{que} o incita a desenvolver a capacidade de pensar problemas filosóficos e criar conceitos.
Marxista, foucaultiano, agostiniano, hegeliano, aristotélico, marcusiano ou orientado por outra corrente de pensamento filosófico, o professor de Filosofia precisa partir do pressuposto de \textbf{que} os objetos de aprendizagem \textbf{que} trabalhará devem proporcionar a formação integral dos \textbf{discentes}.
Essa relação com o saber vai ao encontro do \textbf{que} o filósofo alemão Immanuel Kant entende quando afirma \textbf{que} ensinar Filosofia é diferente de se ensinar a filosofar.
A experiência \textbf{conceitual} \textbf{que} a Filosofia proporciona, entendida por Gallo ( 2012 ) como o objetivo do seu ensino, \textbf{necessita} ser repensada, principalmente no \textbf{que} se refere às metodologias, diante da flexibilização curricular implementada com a lei do Novo Ensino Médio.
Se, na hipótese anteriormente levantada, por semana só haverá 15 horas para a BNCC, para as novas escolas de tempo integral a parte flexível terá, no mínimo, 4 horas aula a mais por dia.
Durante a semana, desse modo, o professor poderá ter \textbf{uma} quantidade a mais de horas aulas para aprofundar o estudo dos conteúdos filosóficos.
Essas horas poderão ser preenchidas por unidades curriculares e eletivas \textbf{que} compõem os itinerários formativos, podendo ser desenvolvidas na forma de laboratórios, oficinas, clubes, observatórios, incubadoras, núcleos de criação artística e núcleo de estudos \textbf{que} aprofundem os objetos de conhecimento presentes na Base, adequado ao contexto de cada instituição escolar.
Neste cenário desafiador, os professores de Filosofia podem criar formas de aprofundar os conhecimentos filosóficos ou testar metodologias especialistas em ensino de Filosofia propõem, assim como as sugestões presentes nos livros didáticos.
Várias dessas metodologias e sugestões podem ser encontradas nas obras de Chauí, Lígia, Gallo, Walter Kohan, Juvenal Savian, Maria Lúcia Graça Aranha, Lipman e outros \textbf{que} sugerem diversas atividades.
Se, no passado, por falta de tempo era difícil a aplicação dessas sugestões \textbf{metodológicas}, em \textbf{uma} realidade na qual horas aulas podem surgir para \textbf{um} aprofundamento de estudos, abre -se novas possibilidades.
Junto a essa prática de experimentação docente, em sintonia com a competência geral nove da BNCC \textbf{que} versa sobre a empatia e a cooperação, deve -se compartilhar as vivências, experiências e pesquisas para \textbf{que} o êxito do ensino de Filosofia na base comum e na parte flexível seja real, valorizando sua presença na educação e possibilitando \textbf{que} se fortaleça como \textbf{um} saber fundamental para a formação integral das alunas/os.
Considerando o caráter “ abstrato ” do conhecimento filosófico ( \textbf{enquanto} construção/reconstrução racional sobre o mundo ) \textbf{que}, para muitos aprendentes, é distante da sua realidade, o professor precisa articular o rigor \textbf{conceitual} da Filosofia, a compreensão dos problemas e conceitos e a busca por instrumentos e meios \textbf{que} tornem significativos e atraentes o aprendizado dos conhecimentos, saberes e práticas do componente Filosofia.
É \textbf{preciso}, para isso, \textbf{que} o professor perceba \textbf{que} a Filosofia é viva e dialoga com a realidade dos estudantes.
O fomento dessa sensibilidade em \textbf{ver} no dia a dia o campo no qual emergem os problemas filosóficos incita \textbf{uma} atenção ao uso de novas estratégias \textbf{metodológicas} \textbf{que} proporcione ao \textbf{discente} vivenciar \textbf{uma} efetiva experiência do pensamento, na qual é catalisada \textbf{um} processo de aprendizagem onde os conteúdos são contextualizados, entendendo \textbf{contextualização} como o ato de vincular o conhecimento à sua origem e a sua aplicação.
Todo esse esforço possibilita \textbf{que} temas e conceitos próprios da Filosofia, como a abstração, possam ser experienciados diferentemente no pensamento.
Sabemos da dificuldade \textbf{que} há, no ensino de Filosofia, em contextualizar problemas e conceitos abstratos, o \textbf{que} dificulta nosso trabalho no Ensino Médio, mesmo daqueles \textbf{que} utilizam práticas inovadoras de ensino. \textbf{Contudo}, torna -se momento importante da aula para \textbf{que} o aprendizado seja potencializado, já \textbf{que} o uso de técnicas e estratégias auxiliam na sua execução.
Neste âmbito, o laboratório, a vivência e a experimentação são recursos pedagógicos \textbf{que} estimulam a integração dos alunas/os às metodologias definidas pelos professores.
Diante dessas dificuldades, o professor de Filosofia pode contorná -las com formação \textbf{que} atenda às necessidades \textbf{exigidas} no ambiente escolar, em torno do uso de recursos tecnológicos, filmes, vídeos, revistas em quadrinhos, músicas, poemas, brinquedos, arte em geral e outros instrumentos \textbf{que} \textbf{explicite} a questão da abstração e leve o \textbf{discente} a se sentir partícipe do processo de \textbf{ensino-aprendizagem}.
Segundo Gallo ( 2012, p. 96 ), isso possibilitaria o \textbf{que} chama de sensibilização, momento da aula “ de chamar a atenção para o tema trabalhado, criar \textbf{uma} empatia com ele, fazer com \textbf{que} o tema ‘ afete ’ os alunas/os ”.
Importante ressaltar \textbf{que} o incentivo a utilização de recursos didáticos \textbf{que} dinamizam as aulas de Filosofia, em diálogo com as aulas expositivas, não busca reduzir o rigor \textbf{conceitual} do conhecimento filosófico.
O \textbf{que} se orienta é a criação de novas relações entre professor e a/o aluna/o \textbf{que} potencialize o processo de aprendizado da Filosofia pelo aprendente do Ensino Médio.
O rigor \textbf{conceitual}, tão característico da Filosofia, necessariamente se realiza lançando mão da história da Filosofia.
A busca de respostas \textbf{que} transcendam a aparência das coisas já aponta para a construção de \textbf{uma} atitude filosófica em relação ao mundo.
A essa atitude filosófica é \textbf{imprescindível} o recurso da história da Filosofia.
Defender seu uso, \textbf{contudo}, não significa \textbf{que} o professor deve organizar seu plano de aula de forma linear, pois a experiência filosófica é tão importante quanto a \textbf{memorização} da história da Filosofia.
Garantir o acesso dos jovens a essa história não consiste na redução da \textbf{disciplina} a \textbf{um} programa centrado, exclusivamente, no estudo linear das construções elaboradas pelos filósofos a cada época.
Partindo do pressuposto de \textbf{que} o exercício da Filosofia está intimamente associado a \textbf{um} determinado \textbf{método} filosófico, é necessário, também, \textbf{que} os jovens possam ter acesso às diversas \textbf{abordagens} construídas pelos filósofos : \textbf{dialético}, fenomenológico, existencialista, racionalista, etc. e, compreendam \textbf{que} cada \textbf{um} corresponde a \textbf{um} determinado modo de pensar, inexistindo, portanto, a Filosofia, mas filosofias.
Essa pluralidade de filosofias garante, também, \textbf{uma} multiplicidade de \textbf{métodos}, seja o da aula tradicional, mas também de \textbf{uma} oficina, clube de leitura, observatório etc.
A forma de \textbf{abordagem} desses modelos pode ser articulada aos temas ou problemas filosóficos \textbf{que} o professor selecione para seu programa de trabalho ( GALLO, 2008, p. 122 – 123 ).
A atitude filosófica consiste em indagar sobre o \textbf{que} nos inquieta, bem como nos posicionar diante do \textbf{que} \textbf{conseguimos} desvelar, descobrir, conhecer a partir das dúvidas e aflições \textbf{que} movem ao questionamento e à reflexão.
Diante das relações desenvolvidas na sociedade \textbf{contemporânea}, seja no campo econômico, político ou social, \textbf{que} novas questões têm sido colocadas pela \textbf{humanidade}?
E pelos jovens \textbf{que} frequentam as escolas públicas do Ceará?
O \textbf{atual} modelo de sociedade tem permitido aos jovens construir atitudes filosóficas em relação ao mundo existente?
Permite pensar e lutar por outras relações possíveis?
Como estimular as/os alunas/os do Ensino Médio à reflexão filosófica no contexto social em \textbf{que} ideias, gostos e comportamentos são fabricados e consumidos em massa?
Introduzir os jovens do Ensino Médio na utilização de procedimentos filosóficos frente a sociedade \textbf{contemporânea} é \textbf{um} dos grandes desafios \textbf{que} se coloca para o componente curricular.
Assim, esses procedimentos filosóficos deverão lançar -se sobre as ciências, a arte, os comportamentos e outras esferas da produção humana, e a entrada nesses procedimentos filosóficos pressupõe o domínio de \textbf{um} referencial \textbf{teórico} e \textbf{metodológico}.
Desse modo, o uso de recursos \textbf{que} potencializam o estudo dos objetos de aprendizagem, assim como sua \textbf{contextualização}, fazem as/os alunas/os adquirir competências e habilidades da Filosofia a partir das suas diversas narrativas.
O docente precisa sentir -se motivado a diversificar o uso, a forma e os espaços de produção do conhecimento das suas alunas/os.
Outra prática \textbf{que}, desde os PCN’s em 1999, vem sendo tema de discussões pedagógicas é o diálogo entre as diversas áreas e seus componentes : a \textbf{interdisciplinaridade}.
A Filosofia, com suas diversas narrativas, possibilita \textbf{um} diálogo no qual estejam relacionados diversos componentes curriculares, sendo \textbf{um} ponto \textbf{que} explana a importância do ensino de Filosofia na educação básica.
A \textbf{interdisciplinaridade}, como \textbf{um} princípio do Ensino Médio, conforme as Diretrizes Curriculares Nacionais ( DCNEM ), tem no estudo e nas narrativas filosóficas \textbf{um} \textbf{forte} aliado no desenvolvimento de \textbf{um} trabalho articulado aos demais componentes curriculares.
Nos momentos do planejamento escolar coletivo, os professores de Filosofia poderão elaborar, com professores de outras \textbf{disciplinas} e áreas do conhecimento, propostas de estudos de temas ou problemas \textbf{que} permitam a utilização dos seus saberes para a compreensão do objeto estudado.
Tal iniciativa possibilita ao \textbf{educando} ampliar sua visão acerca da questão estudada, tendo em vista os elementos \textbf{que} serão fornecidos por cada \textbf{um} dos componentes.
Em \textbf{um} projeto interdisciplinar cujo tema seja “ A Origem do Universo ”, o professor de Filosofia poderia articular com docentes de Química, Física, Biologia, Geografia etc., diversas \textbf{abordagens} \textbf{que} poderiam ser dadas ao tema, considerando os conhecimentos sobre o assunto produzidos por esses saberes.
Além disso, nas aulas de Filosofia, poder -se-ia lançar mão da história da Filosofia, apresentando, por exemplo, as buscas empreendidas pelos pré-socráticos no sentido de definir \textbf{uma} causa material para a origem do universo.
Múltiplas são as possibilidades de trabalhos conjuntos entre as áreas e os componentes.
Entretanto, é importante o entendimento sobre as situações em \textbf{que} determinados temas de estudo devem guardar suas especificidades, evitando articulações forçadas entre \textbf{disciplinas} \textbf{que}, às vezes, podem comprometer o processo de aprendizagem dos aprendentes.
Todo esse esforço proporciona, mais do \textbf{que} nunca, a relação do \textbf{discente} com o saber \textbf{que} bem caracteriza a ideia de Filosofia.
Todas essas considerações acerca do sentido da Filosofia, do uso de recursos \textbf{metodológicos}, de Novo Ensino Médio, de BNCC, de \textbf{contextualização}, de \textbf{interdisciplinaridade}, de história de Filosofia, entre outros, se articulam com \textbf{um} momento essencial no processo de \textbf{ensino-aprendizagem}, \textbf{que} na Filosofia é problema latente, mas \textbf{necessita} de muita atenção : a avaliação.
Por exemplo, a falta de tempo limita o trabalho com as alunas/os, a leitura de textos, a resolução de atividade e, consequentemente, sua correção.
Aplicar \textbf{provas} discursivas torna -se, em muitos momentos, \textbf{uma} \textbf{tarefa} difícil para os professores de Filosofia, pois \textbf{um} docente responsável por 27 turmas têm em média 1100 \textbf{discentes}.
As variantes aumentam as dificuldades na \textbf{tarefa} de avaliar, como o aumento de alunas/os com necessidades especiais \textbf{que} \textbf{necessitam} de atenção especializada e aquelas/es \textbf{que}, em geral, têm dificuldade na compreensão de questões abstratas.
Junto a isso, há \textbf{um} grande número de \textbf{discentes} com baixo grau de proficiência em leitura e escrita, o \textbf{que} faz a avaliação em Filosofia, tema a ser urgentemente pesquisado, \textbf{necessitando} de \textbf{uma} atenção redobrada por parte dos professores do componente.
Neste vácuo é \textbf{que} devemos buscar — além das IES e da SEDUC proporcionarem formação, recursos e incentivos — novas estratégias \textbf{metodológicas}, como acerca da avaliação, \textbf{que} proporcionem aos \textbf{discentes} compreender o problema filosófico, relacionar com sua vida e desenvolver competências, habilidades e objetos conhecimentos.
A questão da avaliação se relaciona, direta e indiretamente, com \textbf{um} elemento importante para o ensino de qualquer saber no Ensino Médio, mas \textbf{que} é \textbf{um} problema recente na Filosofia, se compararmos a outros componentes curriculares, como a história e a geografia : o livro didático.
Desde a primeira edição em \textbf{que} o PNLD ofereceu o livro didático de Filosofia, houve \textbf{uma} intensa produção de obras no mercado editorial, apresentando \textbf{uma} série de atividades, práticas de ensino, exercícios extra sala de aulas etc. \textbf{Contudo}, a análise sobre o livro didático deve considerar outros pontos importantes para \textbf{que} o professor de Filosofia utilize este recurso de forma adequada e da melhor maneira, pois o \textbf{que} ocorre, em muitos casos, é o livro didático ditar a organização do plano de aula do professor de Filosofia durante todo o Ensino Médio, o \textbf{que} não é orientado.
O livro didático é \textbf{um} recurso, como \textbf{um} pincel, \textbf{uma} tv, \textbf{um} filme e outros \textbf{que} potencializam o ensino, mas não deve ser a coluna vertebral de nenhum componente curricular, muito menos da Filosofia.
Os docentes de Filosofia precisam criar \textbf{uma} autonomia para escolher e o trabalhar com obras \textbf{que} possibilitem \textbf{um} trabalho diferenciado, \textbf{uma} sensibilização e \textbf{uma} \textbf{contextualização}, como os \textbf{que} ofereceram recursos áudios visuais, pois os livros didáticos devem conter instrumentos pedagógicos “ como fotos, imagens, músicas e obras de arte \textbf{que} colaborem com o aprendizado dos \textbf{discentes} ” ( SOUSA, p. 265 ), além de “ textos de filósofos em diversas partes de suas obras [ ou ] \textbf{um} espaço específico para se trabalhar textos dos filósofos ” ( Idem, p. 261 ).
Essa sugestão acerca da escolha do livro didático deve ser, também, \textbf{uma} incitação para \textbf{que} os professores desenvolvam seus próprios materiais didáticos, partindo dos problemas das escolas, da vida das alunas/os, da complexidade social, econômica, cultural e política das cidades e do estado do Ceará.
Do mesmo modo \textbf{que} as/os estudantes devem ser estimulados a adquirirem \textbf{uma} postura ativa na sala aula, os professores também devem estimular \textbf{vontade}, \textbf{que} pode ser expressa na produção do próprio material didático.
Todos esses recursos são meios dos professores potencializarem sua \textbf{docência} e \textbf{fomentar}, no aprendente, o desenvolvimento de competências \textbf{que} o forme integralmente.
A última diretriz curricular do Ceará, chamada Escola Aprendente ( 2009 ), apresentava \textbf{uma} discursividade em torno do Ensino de Filosofia alinhada ao contexto histórico em \textbf{que} foi produzido : PCN+, OCN, DCN, do parecer 38/2006 do CNE, \textbf{que} incluía, de forma obrigatória, as \textbf{disciplinas} de Filosofia e Sociologia no currículo do Ensino Médio, e a recém criada, na época, lei nº 11.684/2008, \textbf{que} expandia o ensino dessas \textbf{disciplinas} para todos os três anos do ensino médio.
No presente, as discussões são outras, a maioria decorrente da Lei nº 13.415/2017 ( Novo Ensino Médio ) \textbf{que} flexibiliza o currículo do Ensino Médio e retira a obrigatoriedade de todas as \textbf{disciplinas}, com exceção de português e matemática.
Junto a estas, a BNCC se torna a pauta do momento, principalmente com o retorno do “ velho debate ” das competências e habilidades ( MÔNICA, 2018 ), \textbf{que} na atualidade trazem novas e diversas questões daquelas presentes nas propostas curriculares anteriores.
Vale, neste ponto, questionar se as diretrizes do passado, tanto nacionais como estaduais, precisam ser esquecidas ou \textbf{relegadas} a segundo plano, o \textbf{que} de \textbf{imediato} negamos, pois muitos dos objetivos presentes nestes documentos são anseios da formação em Filosofia no presente, ou sofreram atualizações. \textbf{Vejamos} eles : ⁴ Cálculo hipotético, pois cada rede e escola pode dividir a carga horária referente à BNCC adequada a sua realidade.

% matched lemmas: abordagem, amor, aprendiz, argumentar, atual, basilar, cientista, conceitual, conseguir, contemporâneo, contextualização, contudo, dialético, discente, disciplina, dizer, docência, educar, enquanto, ensino-aprendizagem, equidade, exigir, explicitar, fomentar, forte, homem, humanidade, imediato, implicar, imprescindível, integralidade, interdisciplinaridade, lado, memorização, metodológico, moderno, método, necessitar, preciso, prova, que, relegar, singularidade, sujeito, tarefa, territorialidade, teórico, um, ver, vontade
\end{textsample}
